% SPDX-License-Identifier: Apache-2.0
\section{The Registers a Driver Talks To}
\label{sec:x-ethslots}

An interface is the expensive part and an implementation is the cheap
one, so the interface worth having is the one somebody else is
already paying to maintain.

Zephyr's tree contains \code{drivers/ethernet/eth\_litex\_liteeth.c},
a maintained driver for a soft MAC on an FPGA whose slot size is
\code{0x800}, which is exactly this repository's \code{FRAME\_MAX}.
\code{EthSlots} is that driver's register map in hardware, so that
bringing up Ethernet under Zephyr is a port of a file that already
works rather than a design nobody has driven.

\subsection{Two slots a direction, not a ring}

A descriptor ring was the obvious design and it is the wrong one.
Zephyr's driver model never hands a chain to hardware: it calls
\code{net\_pkt\_read(pkt, buf, len)} to copy a frame out into one
flat buffer before transmitting, and \code{net\_pkt\_write(pkt, buf,
len)} to copy one in after receiving. Both take a flat pointer and a
byte count, so the hardware only ever sees one contiguous region per
frame.

So there are two slots each way, used in turn, each a flat buffer of
2048 bytes. A ring buys queue depth rather than throughput, and if it
is wanted later it is an addition: a descriptor carries the same
three values \code{tx\_slot}, \code{tx\_length} and \code{tx\_start}
already carry.

\subsection{Where the buffers are}

LiteEth's slots are memory inside the peripheral, and its driver
copies each frame into them over the bus. That is the cost this is
meant to remove, only at word granularity rather than byte.

Here the slots are a region of main memory, and the peripheral reads
from it with \code{LineFetch} and fills it with \code{LineStore}
(\S\ref{sec:x-dma}, \S\ref{sec:x-dmaw}). The driver cannot tell the
difference: it is handed a base address either way and writes a frame
there. What changes is whether the processor moves the bytes.

The two directions are separate regions, receive first and transmit
4096 bytes above it. They have to be. A slot number alone would make
transmit slot zero and receive slot zero the same address, and a
frame arriving would land on one waiting to go out: the run below
prints \code{0x41001800} and \code{0x41000800} for the two, which is
the check that they do not meet.

\lstinputlisting[rangebeginprefix=//\ begin\{,rangebeginsuffix=\},%
rangeendprefix=//\ end\{,rangeendsuffix=\},includerangemarker=false,%
linerange=state-state,caption={\code{lib/parts/src/ethslots.rs}: the
registers, as LiteEth lays them out.},%
label={lst:x-ethslots}]{lib/parts/src/ethslots.rs}

\subsection{What the run checks, and why each}

The example does what the driver does, in the order the driver does
it.

\emph{The transmit side has no interrupt.} LiteEth's driver disables
it and polls \code{tx\_ready}, because Zephyr's send may block. So
the run reads \code{tx\_ready} as one while idle and as zero while a
frame is going, which is the whole of the transmit handshake.

\emph{The receive interrupt follows the store, not the frame.}
\code{LineStore} holds its \code{running} high until the write
response comes back from memory, and the interrupt is raised on its
fall. The run asserts no interrupt while the store is in flight and
one after it, so a driver told a frame has arrived can read it
without a barrier. The register block takes the engine's busy line
rather than a completion pulse for that reason: a pulse could be
joined to the receiver by somebody reading the map and not the
comment.

\emph{The pending bit is write-one-to-clear, and that is tested both
ways.} Writing zero must not acknowledge, and writing one must.
Nothing about a register's value says which behaviour it has, and
hardware that cleared on any write would let a driver acknowledge an
interrupt it had not looked at, dropping the frame rather than
reading it. Only a run that writes zero and finds the bit still set
can tell those apart.

Four signals say the whole of it: \code{tx\_go} while the transmit
request waits for the engine to take it, \code{rx\_pending} from the
arrival to the acknowledgement, and \code{irq} following it.

\input{ethslots_timing.tex}

\subsection{The cycle where the two meet}

The registers are written with \code{with!}, whose entries apply in
order with the last drive of a field winning. That makes the order of
two of them a decision about frames rather than about style.

An arrival sets \code{rx\_pending} and an acknowledgement clears it.
Written in the order they read most naturally, arrival first, an
acknowledgement in the same cycle as an arrival clears the bit the
arrival has just set. The frame is in memory, whole and correct, and
nothing ever tells the driver it is there. It is the cycle where a
busy link acknowledges one frame as the next one lands, so the fault
appears under load and not in a quiet test.

So the acknowledgements are written first and the arrival after them,
and a frame arriving as one is acknowledged survives.

The test for it was vacuous twice before it was worth anything, which
is the part worth reporting. An acknowledgement aimed at the arrival
by reasoning about the bus latency landed a cycle early, where the
arrival sets the bit afterwards whichever order the entries are in.
The example passed against the faulty order, which is the only reason
the mistake was visible at all.

What settled it was running the whole sweep against both orders:

\begin{center}\footnotesize
\begin{tabular}{@{}lccccc@{}}
\toprule
issued at & 199 & 200 & 201 & 202 & 203 \\
\midrule
acknowledgements first & ok & ok & ok & fails & fails \\
arrival first & ok & ok & fails & fails & fails \\
\bottomrule
\end{tabular}
\end{center}

At 199 and 200 the acknowledgement is early and proves nothing. At
202 and later it falls in a cycle of its own, after the arrival,
where clearing the bit is what should happen and both orders fail.
One cycle out of the five, 201, is the one where the two coincide,
and it is the only column that tells the orders apart.

A delay chosen so that a test passes is worth nothing, and this one
is chosen so that it fails on the fault it names.

\subsection{What the driver must still do}

Stores on this machine are posted and \code{fence} orders nothing
(issue 432), so a driver must read the last word of a frame back
before writing \code{tx\_start}. Otherwise the engine may fetch
before the frame has landed and send what was in the slot before.

That obligation is on the component's datasheet rather than only in
the driver, because the next driver written against this peripheral
will not read the first one. What makes it dischargeable is measured
in the AXI document~\cite{axi}: this link answers a load behind an
unanswered store to the same address with the stored value, which AXI
does not require and this interconnect does.

\lstinputlisting[style=ascii,caption={What \code{ex\_ethslots}
printed when the build ran it.},%
label={lst:x-ethslots-out}]{out_ethslots.txt}
